% Suggested final budget: 1.2 pages. Merges the former Method and Experimental
% Design sections; keep each subsection to one or two short paragraphs.
%
% Drafted 2026-10-12 from the repository, not from run records. Sources:
% dp-manip/configs/baseline.toml and configs/experiments/*.toml (hyperparameters,
% seeds), dp-manip/dp_manip/{policy,vision,observation_encoder,data}.py and
% backbones/*.py (architecture, normalization, sampling), dp-manip/PLAN.md
% (protocol), configs/failure_aware/takeover_*.toml and
% docs/failure-expert-takeover-probe.zh-CN.md sections 7-9 (takeover study).
% Check each statement against run.json of the evaluated runs.
\section{Method and Experimental Setup}
\label{sec:method}
\subsection{Conditional action denoising}
For observation history $o$ and clean action sequence $a^{(0)}$, the policy is
trained with the noise-prediction objective of Diffusion
Policy~\citep{chi2023diffusion,ho2020ddpm},
\begin{equation}
  \mathcal{L}(\theta)=
  \mathbb{E}_{o,a^{(0)},k,\epsilon}
  \left[\left\|\epsilon-
  \epsilon_\theta\!\left(
  \sqrt{\bar\alpha_k}a^{(0)}+
  \sqrt{1-\bar\alpha_k}\epsilon,k,o\right)
  \right\|_2^2\right],
  \label{eq:denoising}
\end{equation}
where $k$ is the diffusion timestep and $\epsilon\sim\mathcal{N}(0,I)$. The
policy observes the last $T_o=2$ steps, predicts $T_p=16$ actions, and executes
$T_a=8$ of them, starting at the current step, before replanning. Each action
dimension is scaled to $[-1,1]$ with the minimum and maximum over the training
subset and mapped back before execution. We use the squared-cosine
noise schedule~\citep{nichol2021improved} with 100 steps for both training and
sampling.

\subsection{Observation encoder and backbones}
Each RGB frame is scaled to $[-1,1]$ and encoded by a
ResNet-18~\citep{he2016resnet} with group normalization~\citep{wu2018groupnorm},
trained from scratch, followed by global average pooling and a projection to
128 features; cameras share one encoder. Images are randomly shifted by up to
4 pixels during training~\citep{kostrikov2021drq}. The image features are concatenated
with the standardized low-dimensional state, giving one feature vector per
observation step.

The three noise-prediction backbones receive the same feature sequence and
differ only in how they consume it. The \textbf{UNet} is the 1D convolutional
network of Diffusion Policy as shipped with the ManiSkill baseline, used
unmodified (channel widths 256/512/1024, kernel size 5), with
FiLM conditioning~\citep{perez2018film} on the flattened observation
features. The \textbf{Transformer} is a causal
decoder over action tokens that attends to a timestep token and the observation
tokens (8 layers, 4 heads, width 256, attention dropout 0.3). The \textbf{MLP}
has three hidden layers of width 256 over the concatenation of the flattened
noisy actions, a timestep embedding, and a 256-dimensional projection of the
observation features.
The Transformer follows the transformer variant of the Diffusion Policy
repository and the MLP is our own baseline design; both were ported from an
earlier codebase of our group without changing their structure.

\subsection{Shared training configuration and compute}
All runs share one configuration: AdamW~\citep{loshchilov2019adamw} with
learning rate $10^{-4}$,
$\beta=(0.95,0.999)$ and weight decay $10^{-6}$; 500 warm-up steps followed by
cosine decay; batch size 64; gradient clipping at 1.0; mixed precision; and an
exponential moving average of the weights with decay 0.9999. Every run trains
for 100k optimizer steps regardless of task, backbone, or $N$, and we evaluate
the averaged weights of the final checkpoint. Checkpoints at 10k, 30k, and 60k
steps are kept for diagnostics only. The only per-task settings are the control
mode and the episode horizon (\cref{tab:tasks}).

Each run trains on a single GPU. On an RTX~4080 SUPER a 100k-step UNet run
takes 1.7~h and peaks at 2.6~GB of GPU memory (\cref{sec:backbone}). The
data-size and backbone studies comprise 166 such runs; as one example of the
per-task budget, the 26 training runs of PushCube and their evaluations took
48.7 GPU-hours.

\subsection{Evaluation protocol}
Policies are evaluated in closed loop on the same CPU physics backend that
generated the demonstrations. Each trained run is tested on 100 episodes with
reset seeds 10000--10099, held out as described in \cref{sec:separation}.
Sampling noise uses a fixed seed.
\metriconce{} counts an episode in which the success condition holds at any
step and is our primary metric; \metricend{} requires it to hold at the final
step. Validation loss and rollouts from training seeds serve only as
diagnostics.

We report the mean and sample standard deviation over training seeds.
Differences between adjacent data sizes or between backbones are assessed with
a two-level bootstrap that resamples training seeds and then evaluation
episodes; we call a difference distinguishable when its 95\% interval excludes
zero.
Exceptions to this procedure and to the shared horizon and split are listed
in \cref{sec:limits}.

\subsection{Study designs}
\textbf{Data size (RQ1).} For each task we train the UNet baseline on
$N\in\{25,50,100,200\}$ demonstrations, with 3 training seeds for $N\le50$ and
5 for $N\ge100$. Subsets are nested and fixed across training seeds
(\cref{sec:separation}), so the training seed changes only the initialization
and batch order. $N=400$ (5 seeds) is added for a
task only when its $100\rightarrow200$ gain is distinguishable from zero, a
rule fixed before the runs. Because the step budget is fixed, smaller subsets
are revisited more often rather than trained for fewer steps.

\textbf{Backbone (RQ2).} At $N=100$ we compare the three backbones with 5
training seeds each. The UNet arm reuses the data-size runs, and the resolved
configurations of the three arms differ only in the backbone selector. The
backbones differ widely in size (\cref{sec:backbone}), so the comparison is
between these implementations, not between architecture families at equal
capacity.
The comparison covers PickCube, StackCube, PushCube, PullCube, and
LiftPegUpright; it was not run on PlaceSphere.

\textbf{Takeover correction (RQ3).} On PlaceSphere, starting from the three
$N=100$ UNet baseline checkpoints, a trigger that reads privileged simulator
state detects a missed grasp during a rollout (the gripper has closed on
nothing and the sphere has not been lifted), and a scripted expert takes over,
regrasps, and places the sphere. For each checkpoint we collect 100 successful
corrections from its own rollouts. Three arms are compared: B, the baseline;
C, fine-tuned on the original 100 demonstrations plus the corrections; and D,
fine-tuned on the original 100 plus ordinary demonstrations with at least as
many training windows. C and D share the fine-tuning settings (learning rate
$3\times10^{-5}$, 10k steps), and all arms are evaluated on the same 200 new
seeds per checkpoint. Trigger, expert, arms, and decision rule were fixed
before the corrections were collected.
